\[ %%%%%%%%%%%%%%%%%%%%%%%%%%%%%%%%%%%%%%%%%%%%%%%%%%%%%%%%%%%%%%%%%%%%%%%%%%%%%%%% \subsection{Computational Complexity and Scalability} \label{sec:complexity} %%%%%%%%%%%%%%%%%%%%%%%%%%%%%%%%%%%%%%%%%%%%%%%%%%%%%%%%%%%%%%%%%%%%%%%%%%%%%%%% The scalability of physics-guided learning architectures is a critical factor for their application to large-scale scientific simulations and digital twin systems. Physical systems of practical interest often contain millions of interacting components, requiring computational frameworks capable of efficiently processing large dynamic structures. PHYSGEN is designed to preserve the computational advantages of graph-based learning while introducing additional physics-guided operations required for physically consistent prediction. %%%%%%%%%%%%%%%%%%%%%%%%%%%%%%%%%%%%%%%%%%%%%%%%%%%%%%%%%%%%%%%%%%%%%%%%%%%%%%%% \subsubsection{Computational Complexity of Graph Evolution} %%%%%%%%%%%%%%%%%%%%%%%%%%%%%%%%%%%%%%%%%%%%%%%%%%%%%%%%%%%%%%%%%%%%%%%%%%%%%%%% Let a physical system be represented by a graph: \begin{equation} G=(V,E), \end{equation} where $N=|V|$ represents the number of nodes and $M=|E|$ represents the number of edges. The computational complexity of standard message passing in graph neural networks is approximately: \begin{equation} \mathcal{O}(M d), \end{equation} where $d$ denotes the latent feature dimension. Since physical interaction graphs are typically sparse, where: \begin{equation} M \ll N^2, \end{equation} graph-based approaches provide significantly better scalability compared with dense attention mechanisms. The core PHYSGEN message passing operation maintains this sparse complexity: \begin{equation} \mathcal{O}_{msg} = \mathcal{O}(M d). \end{equation} %%%%%%%%%%%%%%%%%%%%%%%%%%%%%%%%%%%%%%%%%%%%%%%%%%%%%%%%%%%%%%%%%%%%%%%%%%%%%%%% \subsubsection{Additional Cost of Physics Guidance} %%%%%%%%%%%%%%%%%%%%%%%%%%%%%%%%%%%%%%%%%%%%%%%%%%%%%%%%%%%%%%%%%%%%%%%%%%%%%%%% The introduction of physics-guided modulation adds additional computational operations. The physics interaction coefficient: \begin{equation} \alpha_{ij}^{phys} = \Psi_{phys} (h_i,h_j,e_{ij},P_{ij}) \end{equation} requires evaluation for each graph edge. Therefore, the complexity of physics guidance is: \begin{equation} \mathcal{O}_{phys} = \mathcal{O}(M d_p), \end{equation} where $d_p$ represents the dimensionality of physical descriptors. The total complexity of a PHYSGEN layer can therefore be approximated as: \begin{equation} \mathcal{O}_{PHYSGEN} = \mathcal{O} ( M(d+d_p) ). \end{equation} Thus, physics guidance introduces an additional linear cost with respect to the number of physical interactions. %%%%%%%%%%%%%%%%%%%%%%%%%%%%%%%%%%%%%%%%%%%%%%%%%%%%%%%%%%%%%%%%%%%%%%%%%%%%%%%% \subsubsection{Complexity of Latent Dynamics Evolution} %%%%%%%%%%%%%%%%%%%%%%%%%%%%%%%%%%%%%%%%%%%%%%%%%%%%%%%%%%%%%%%%%%%%%%%%%%%%%%%% The latent dynamics operator: \begin{equation} H_{t+1} = \Psi_{\theta}(H_t,G_t,P_t) \end{equation} operates on node representations after physics-guided aggregation. For $L$ message-passing layers, the approximate computational complexity becomes: \begin{equation} \mathcal{O} ( L M d ). \end{equation} The complexity remains linear with respect to graph size, enabling application to large sparse physical systems. %%%%%%%%%%%%%%%%%%%%%%%%%%%%%%%%%%%%%%%%%%%%%%%%%%%%%%%%%%%%%%%%%%%%%%%%%%%%%%%% \subsubsection{Long-Horizon Prediction Cost} %%%%%%%%%%%%%%%%%%%%%%%%%%%%%%%%%%%%%%%%%%%%%%%%%%%%%%%%%%%%%%%%%%%%%%%%%%%%%%%% Long-horizon forecasting requires recursive application of the evolution operator. For prediction horizon $K$: \begin{equation} H_{t+K} = \Psi_{\theta}^{K}(H_t), \end{equation} the computational cost becomes: \begin{equation} \mathcal{O}_{forecast} = K \mathcal{O}_{PHYSGEN}. \end{equation} However, the adaptive stability controller reduces unnecessary error correction by dynamically regulating the evolution process. Future extensions may incorporate adaptive temporal stepping: \begin{equation} \Delta t=\eta(s_t), \end{equation} allowing computational resources to be concentrated on dynamically complex regions. %%%%%%%%%%%%%%%%%%%%%%%%%%%%%%%%%%%%%%%%%%%%%%%%%%%%%%%%%%%%%%%%%%%%%%%%%%%%%%%% \subsubsection{Scalability Strategies} %%%%%%%%%%%%%%%%%%%%%%%%%%%%%%%%%%%%%%%%%%%%%%%%%%%%%%%%%%%%%%%%%%%%%%%%%%%%%%%% PHYSGEN supports several strategies for large-scale deployment: \paragraph{Sparse Graph Computation.} Physical interactions are naturally local in many domains. Sparse adjacency representations allow efficient memory usage and parallel processing. \paragraph{Hierarchical Graph Representation.} Large systems can be decomposed into multiple spatial or functional scales: \begin{equation} G= \{G_1,G_2,\dots,G_k\}, \end{equation} where each subgraph represents a different physical resolution. \paragraph{Distributed Training.} The graph structure enables partition-based parallelization: \begin{equation} G = \bigcup_i G_i, \end{equation} allowing training across multiple computational devices. \paragraph{Adaptive Resolution.} PHYSGEN can allocate computational complexity according to physical importance: \begin{equation} C_i=f_{adaptive}(x_i,s_i), \end{equation} where $C_i$ represents local computational effort. %%%%%%%%%%%%%%%%%%%%%%%%%%%%%%%%%%%%%%%%%%%%%%%%%%%%%%%%%%%%%%%%%%%%%%%%%%%%%%%% \subsubsection{Comparison with Existing Approaches} %%%%%%%%%%%%%%%%%%%%%%%%%%%%%%%%%%%%%%%%%%%%%%%%%%%%%%%%%%%%%%%%%%%%%%%%%%%%%%%% The computational characteristics of PHYSGEN can be summarized relative to existing scientific machine learning approaches. \begin{table}[h] \centering \caption{Computational characteristics of scientific learning approaches.} \label{tab:complexity_comparison} \begin{tabular}{lccc} \hline \textbf{Method} & \textbf{Physics Integration} & \textbf{Scaling Behavior} & \textbf{Long-Horizon Capability} \\ \hline PINN & High & Limited by optimization & Limited \\ Dense Transformer & Low/Optional & $\mathcal{O}(N^2)$ & Medium \\ Neural Operator & Medium & Problem dependent & Medium \\ Graph Neural Simulator & Low/Medium & $\mathcal{O}(M)$ & Medium \\ \textbf{PHYSGEN} & High & $\mathcal{O}(M(d+d_p))$ & Designed for long horizon \\ \hline \end{tabular} \end{table} %%%%%%%%%%%%%%%%%%%%%%%%%%%%%%%%%%%%%%%%%%%%%%%%%%%%%%%%%%%%%%%%%%%%%%%%%%%%%%%% \subsubsection{Summary} %%%%%%%%%%%%%%%%%%%%%%%%%%%%%%%%%%%%%%%%%%%%%%%%%%%%%%%%%%%%%%%%%%%%%%%%%%%%%%%% The computational design of PHYSGEN balances physical consistency with scalability. By maintaining sparse graph operations and integrating physics guidance through local interaction mechanisms, the framework avoids the computational limitations associated with dense global architectures. The resulting architecture provides a scalable foundation for modeling large nonlinear dynamical systems, including scientific simulations, engineering platforms, and physics-based digital twins. \] 